\section{Proposed method: directional group-peak updates}\label{sec:directional}
The second construction couples interval departures within disjoint
groups. A shared peak price can allocate corrections differently from
independent row prices: changing the largest weighted departure alters
the group cost, and several rows may become active at the same peak.
Unequal upper and lower scales specify directional policy. The interface
retains original interval rows while using an exact group proximal
operation. Directional expressibility is not new: a one-sided row split
or a compact epigraph already expresses the same objective. The research
question concerns the resulting complete numerical operation.

For a nonempty closed box $C$ and a residual penalty $\Psi$, define
\begin{equation}\label{eq:paper-infimal}
 g_C(z)=\inf_{s\in C}\Psi(z-s).
\end{equation}
Assume $\Psi$ is proper, closed, convex, nonnegative and satisfies
$\Psi(0)=0$. The relevant contraction property is
\begin{equation}\label{eq:paper-contraction}
 \Psi(Dw)\leq\Psi(w),\qquad D=\diag(d_i),\quad0\leq d_i\leq1.
\end{equation}
This property permits unequal behavior on positive and negative rays.
It does not require global sign symmetry. Orthant-monotonicity provides
important prior context for that distinction~\cite{orthant}; convexity
alone is not enough to justify the box projection identity below.

\begin{theorem}[Box-residual identity]\label{thm:paper-box}
The contraction property\eqref{eq:paper-contraction} holds if and only if,
for every nonempty closed box and every $z$,
\begin{equation}\label{eq:paper-box-identity}
 g_C(z)=\Psi(z-\Pi_C(z)).
\end{equation}
\end{theorem}
\begin{proof}
For any $s\in C$, set $r=z-\Pi_C(z)$. If $z_i$ exceeds its upper bound,
then $z_i-s_i\geq r_i>0$. If it is below its lower bound, then
$z_i-s_i\leq r_i<0$. Inside the interval $r_i=0$. Therefore
$r=D(z-s)$ for some diagonal contraction. The property gives
$\Psi(r)\leq\Psi(z-s)$, with equality attained by $s=\Pi_C(z)$.
Conversely, fix $w,D$ and choose coordinate intervals with endpoints
$-w_i$ and $-d_iw_i$ in increasing order. At $z=0$ the nearest-box
residual is $Dw$, while $-w$ is an admissible point. The claimed identity
therefore implies $\Psi(Dw)\leq\Psi(w)$.
\end{proof}

The theorem states an equivalence over all boxes. It is not merely a
numerical observation on one training case. The infimal-convolution
representation gives convexity. Its explicit composition with the
continuous residual map supplies lower semicontinuity under the stated
closedness assumption. A point inside the box has zero price, so the
eliminated penalty is proper. These domain statements matter when an
extended-real penalty represents an immutable restriction rather than
a finite departure price.

Let $H=\diag(h_i)\succ0$ be fixed. Define
\begin{equation}\label{eq:paper-metric-prox}
 \prox_f^H(v)=\argmin_z\{f(z)+\tfrac12(z-v)^\top H(z-v)\}.
\end{equation}
The translated residual relation is
\begin{equation}\label{eq:paper-translated-prox}
 \prox_{g_C}^H(v)=\Pi_C(v)+\prox_{\Psi}^H(v-\Pi_C(v)).
\end{equation}
\begin{proof}
Put $p=\Pi_C(v)$ and clip a candidate coordinate to the
segment joining $p_i$ and $v_i$. Diagonal quadratic distance from $v$
cannot increase under that clipping. The new box residual is a
coordinate contraction of the original residual, so its penalty cannot
increase either. An optimum can therefore be sought on that product
of segments. Writing $z=p+w$ makes its residual exactly $w$. The
unconstrained residual proximal minimizer lies on the same segments
because clipping toward the segment from zero to $v-p$ again decreases
both terms. Thus the segment restriction does not change the residual
minimum, establishing the relation.
\end{proof}

Diagonal structure is essential to this proof. A dense metric couples
coordinates, and coordinate clipping can then increase the quadratic
distance. The theory section gives a concrete counterexample. Likewise,
the identity is not claimed for an arbitrary convex price lacking the
contraction property. The admitted directional group penalties satisfy
it because reducing a coordinate magnitude without changing its sign
cannot increase any scaled directional peak. Separate disjoint groups
then permit independent group calculations within the diagonal metric.

For one group, define
\begin{equation}\label{eq:paper-directional-price}
 M(w)=\max_i\{c_{U,i}(w_i)_+,c_{L,i}(-w_i)_+\},\qquad
 \Psi(w)=aM(w)+\tfrac12bM(w)^2,
\end{equation}
with positive directional scales and $a,b\geq0$. For residual $r=v-p$,
choose $c_i$ according to the sign of $r_i$. Fixing a candidate peak
$t\geq0$ constrains each residual coordinate to its corresponding
interval. The best coordinate and scalar objective are
\begin{equation}\label{eq:paper-threshold-model}
 w_i(t)=\operatorname{sign}(r_i)\min\{|r_i|,t/c_i\},\qquad
 F(t)=at+\tfrac12bt^2+\tfrac12\sum_i h_i(|r_i|-t/c_i)_+^2.
\end{equation}
The remaining problem is convex in one scalar. Breakpoints occur at
$c_i|r_i|$. Sorting them in descending order gives active prefixes. On an interval with
active set $S(t)=\{i:c_i|r_i|>t\}$, differentiating yields
\begin{equation}\label{eq:paper-threshold-root}
 t=\frac{\sum_{i\in S(t)}h_i|r_i|/c_i-a}
          {b+\sum_{i\in S(t)}h_i/c_i^2}.
\end{equation}
Only a root consistent with its active interval is accepted. The zero
endpoint is checked explicitly. When both prices vanish the original
residual is returned directly; an empty active set never triggers a
division by zero. Repeated breakpoints and zero residuals require
consistent endpoint handling rather than arbitrary active-set choices.

The operation costs $O(k\log k)$ for sorting a group of size $k$, then
a prefix scan. Weighted threshold and projection methods already employ
such structures~\cite{weighted,duchi,condat}. The implementation's specific
role is to preserve the directional interval meaning, diagonal metric
and original-row split, including boundary and representability cases.
No claim identifies sorting or the scalar threshold primitive itself
as an original algorithmic discovery.

Algorithm~\ref{alg:paper-directional} makes those cases explicit. Numerical
threshold comparisons use finite local spacing rather than a fixed
absolute departure floor that changes meaning under unit conversion.
Nonrepresentable arithmetic is refused. The interface requires disjoint
groups; treating overlaps as separate independent calls would in general
change the objective's proximal operation. Positive finite scales are
required, and zero group price is distinguished from a hard constraint.
These are semantic conditions of the implemented method rather than
details to infer from a successful backend return code.

\begin{algorithm}[htbp]
\caption{Evaluate the directional disjoint-group proximal map with a fixed diagonal metric, retaining exact endpoint, zero-price and original-row semantics for independent numerical acceptance.}\label{alg:paper-directional}
\begin{algorithmic}[1]
\Require Nonempty box, disjoint groups, positive scales and diagonal $H\succ0$
\State Set $p=\Pi_C(v)$ and $r=v-p$
\For{each group}
\If{both group prices are zero}
\State Keep that residual group unchanged
\Else
\State Select each scale from the residual sign; sort $c_i|r_i|$
\State Check $t=0$ and consistent roots of \eqref{eq:paper-threshold-root}
\State Recover $w$ using \eqref{eq:paper-threshold-model}
\EndIf
\EndFor
\State Return $p+w$, retaining exact hard-row projection separately
\end{algorithmic}
\end{algorithm}

The joint conjugate is needed for independent acceptance. At fixed peak,
maximizing the dual pairing produces the directional quantity
\begin{equation}\label{eq:paper-dual-budget}
 L(y)=\sum_i\left((y_i)_+/c_{U,i}+(-y_i)_+/c_{L,i}\right).
\end{equation}
Taking the scalar conjugate gives
\begin{equation}\label{eq:paper-peak-conjugate}
 \Psi^*(y)=\begin{cases}
 (L(y)-a)_+^2/(2b),&b>0,\\
 0,&b=0,\ L(y)\leq a,\\
 +\infty,&b=0,\ L(y)>a.
 \end{cases}
\end{equation}
The interval support term remains necessary: $g_C^*=\sigma_C+\Psi^*$.
A single coordinate-wise price cap is therefore not a valid replacement
for the group budget. The evaluator checks each group's absolute
normalized Fenchel relation, raw and bounded-repaired stationarity,
and every hard row's complementarity separately. A negative approximate
hard-row gap cannot cancel an incorrect soft relation.

Finally, exact competing formulations delimit the computational claim.
Directional one-sided rows reproduce the same peak under the published
grouped projection path, while one nonnegative epigraph variable per
group represents the convex price in an ordinary QP. The registered
comparison retains both alternatives and direct/proximal DAQP policies
where compatible. The primary operation includes 32 updates and all
required recovery/acceptance work, not only one scalar threshold call.
Its negative complete-population score demonstrates that a smaller
original-row system can coexist with worse measured complete-operation
cost. It does not refute the proximal identity, imply that every
instance is slower by the same factor, or establish an advantage on
an untested application regime.

Absent interval sides require explicit domain treatment. A missing upper
bound removes its positive directional departure, and a missing lower
bound removes its negative departure. It does not contribute an artificial
large finite endpoint or a product involving infinity and zero. The
interval support in the conjugate then restricts the corresponding dual
sign. This links a modeling convention at the serialized input boundary
to a mathematical condition in original-space acceptance. A primal answer
with a plausible objective can still fail if its recovered multiplier
assigns force to a side that does not exist.

Free group prices provide another useful boundary case. If both the
linear and quadratic group coefficients vanish, the residual penalty
is identically zero and the unconstrained residual proximal map is the
identity. This is different from a hard box indicator, whose row operation
projects to the box. The conjugate of the free residual penalty permits
only zero soft-group dual contribution. These relationships give mutually
consistent primal, proximal and dual controls. They prevent an endpoint
special case from accidentally converting a free preference into an
immutable constraint.

Within a nonzero-price group, the scalar threshold can make several
coordinates meet the same weighted peak. A coordinate below the active
threshold keeps its residual value, whereas a larger weighted residual
is clipped to its directional peak allowance. This is a coordinated
allocation, not an independent softening of each original row. At tied
breakpoints several algebraic active prefixes can describe the same
unique proximal vector. The implementation must recover that vector
consistently rather than infer scientific nonuniqueness from an internal
sorting tie. Empty groups and overlapping memberships are rejected by
the actual interface.

The weights in\eqref{eq:paper-threshold-root} have physical meaning.
The metric contribution divides by the squared directional scale,
while the linear numerator divides by that scale once. Omitting one
factor changes both threshold location and dimensional covariance.
The independently expressed epigraph gives a different numerical route
to the same group optimum, making this a concrete falsification target.
The physical teaching input compares actual native, row-split and
epigraph outputs and retains the original hard coupling equality.
It is deliberately small enough to inspect without treating agreement
as additional evidence of timing superiority.

The original-row formulation also does not imply constant preparation
cost under every update. A change in a vector is different from a
change in matrix structure, metric, group membership or penalty family.
Only the registered repeated-forecast vectors are updated within a
prepared operation. Hard dynamics and heater-energy limits remain
immutable in the thermal family. A different group layout or unit
mapping would identify another input batch and invalidate reuse of
the original source-bound observation. This is why batch digests,
units and method policy settings accompany each recorded operation.

Compatibility of references is similarly derived from semantics, not
from an optimistic output flag. The linear row-split reference cannot
be used as if it handled every quadratic group price. A direct DAQP
epigraph requires its admitted curvature conditions; documented
proximal policies have separately fixed roles. The complete protocol
retains excluded cells as incompatible and dependent oracle limits
as oracle-failed. Consequently the fastest selected reference is
defined by the registered compatible alternatives for that instance,
and cannot be chosen from an arbitrary easier objective afterward.

All those conditions concern the same optimization model from
different directions: input units, convex proximal structure,
equivalent formulations, original duals and complete-operation
accounting. Their agreement establishes a restricted and reproducible
implementation claim. It does not authorize transferring a favorable
teaching example to the entire held-out timing population. The retained
negative score includes all required members, including those whose
candidate or reference operation did not yield an accepted complete
answer. That distinction preserves the scientific question even when
a local closed-form calculation is exact.
